\section{Critical paths and actual improvements}\label{sec:critical-paths}

The preceding extension theorem permits a choice of exact problems in which
the relative density exponent and the optimal time cost evolve linearly.  We construct
these paths first.  We then pass to short portions of them, keeping an
explicit exclusion in the original resolution.  The last part of the section
shows how local candidates return to that original exclusion with sufficient
mass, geometric bounds and terminal label knowledge.

\subsection{Extremal exact problems}

Fix the model version and \(j,S,\eta\), and suppose bad problems exist.
We optimize in three stages: approach the least density cap that permits
positive horizon, maximize the horizon relative to length there, and,
in version 2, optimize the narrowness available on the resulting
sequences. The first two stages are expressed by
\begin{equation}
d_*=\inf\{d<1:\exists E\in C(d),\ H(E)>0\},\qquad
 k(d)=\sup_{E\in C(d)}H(E)/L,\qquad
 k=\lim_{d\downarrow d_*} k(d).\label{a03:extremal-rates}
\end{equation}
Here \(0\le d_*<1\); in the right-hand limit the classes are nonempty
and nested by monotonicity.

The supremum defining \(k(d)\) uses all lengths, equivalently length
one. Replacing the base \(N_0\) by \(N_0^L\) divides depth, horizon
and log-density exponents by \(L\), preserving \(C(d)\), including
its angular parameters, and the chart rules. This scaling takes place
within one exact problem, with its positive length \(L\) fixed.


\begin{lemma}[Composition of horizon costs]\label{a03:composition}
Let \(E\in C(d)\) have length \(L\).  If a known atlas at depth \(r<L\)
has horizon exponent \(h\), and \(F\) is a homogeneous remaining problem
in the original exponent units, then
\begin{equation}
H(E)\le h+H(F).
 \label{a03:known-composition}
\end{equation}
If \(H_{\mathrm{out}}\) is the horizon infimum for the coarsening of \(E\)
to endpoint \(r>0\), then
\begin{equation}
H(E)\le H_{\mathrm{out}}+k(d)(L-r).
 \label{a03:coarsened-composition}
\end{equation}
These statements allow all the subsequences and incidence restrictions
in the definition of the respective infima.
\end{lemma}
\begin{proof}
Choose a true terminal system in \(F\) with horizon within an arbitrary
positive tolerance of \(H(F)\), on a subsequence where that system exists.
Lift it through the known atlas using the terminal equality in
\Cref{lem:normalization}.  The time exponents add, and the lifted system is
true at the original terminal depth.  Letting the tolerance tend to zero
proves \Cref{a03:known-composition}.

Coarsening to \(r\) keeps the trajectory and incidence data and the profile
on \([0,r]\).  The angular coarsening in
\Cref{a01:angular-coarsening} gives the required angular cap at \(p_r\).  Thus only
the terminal profile-slope condition at \(r\) has to be checked when a cap
class is asserted for the coarsening.  A true terminal system of that
coarsening can be prepared with its label known by \(p_r\); it is then a
true intermediate system for \(E\), and its label is also known by the
finer \(p_L\), with the permitted list loss.

Take these outer true systems with costs tending to \(H_{\mathrm{out}}\).
Normalize the original remaining problem through them.  It belongs to
\(C(d)\), has length \(L-r\), and hence has horizon at most
\(k(d)(L-r)\).  Apply the first inequality and let the outer tolerance
tend to zero.  A restriction to subpower incidence mass leaves the
homogeneous density exponents unchanged.  Moreover, systems on such
restrictions and subsequences are already among the competitors defining
\(H(E)\).  This proves the inequalities with their stated quantifiers.
\end{proof}

\begin{lemma}[A positive critical time rate]\label{a03:positive-rate}
The rates in \Cref{a03:extremal-rates} satisfy \(0<k\le1\).
\end{lemma}
\begin{proof}
The upper bound follows from \Cref{a01:horizon-bound}.
For a bad problem \(E\), choose \(e>0\) small compared with \(H(E)>0\)
and its length, and choose a true terminal system with horizon
\(h<H(E)+e/4\).  By \Cref{lem:terminal-extension}, it gives a known
atlas with horizon \(h-e\) and positive remaining length at most \(Ce\).
Normalize through this atlas.  The remaining problem belongs to \(C(d)\)
and, by \Cref{a03:known-composition}, its horizon is at least
\begin{equation}
H(E)-(h-e)>3e/4.
\end{equation}
Its horizon-to-length ratio is therefore at least \(3/(4C)\).
For \(d\) near \(d_*\), choose one \(d_0<1\) with \(d\le d_0\);
the extension constant is uniform for this range and for the fixed
\(j,S,\eta\).  Bad problems exist arbitrarily close to the defining
dimension infimum.  The lower bound for \(k\) follows.
\end{proof}

Scale units to \(L=1\) when discussing maximizing sequences.  These
are sequences of \emph{exact problems} \(E_i\in C(d_i)\) with
\(d_i\to d_*\) and \(H(E_i)\to k\); they exist by the definition.
We have \(\limsup_i k(d_i)\le k\), by comparison with dimensions
slightly larger than \(d_*\).  For each fixed \(i\), all exponents
are first taken along its exact \(N_0\)-sequence; the outer limit
in \(i\) comes afterward.  A single \((i,N_0)\)-diagonal does not
replace these inner limits.

In version 2, fix one maximizing sequence.  Consider true terminal
systems whose horizon exponents tend to \(k\).  For each such system
within \(E_i\), pass to a homogeneous subfamily of charts to obtain
a narrowness exponent, as in \Cref{a01:narrowness}, in length-one
units.  Define \(\omega\) to be the supremum of the attainable
\(\limsup\) values of these exponents.  The choices include systems
on any subsequence of the outer indices \(i\), and all permitted
subsequences and selections within the individual exact problems.
Choices near the horizon infimum with some homogeneous narrowness
exponent exist, and narrowness is nonnegative.

Put \(\ell=\inf\omega\), where the infimum ranges over all maximizing
sequences and may be infinite.  In version 1 we use \(\ell=0\)
formally, without optimizing width.

\begin{proposition}[Critical paths]\label{prop:critical-path}
Fix the model version and its parameters \(j,S,\eta\), and suppose bad
problems exist.  With \(d_*,k,\ell\) as above, the following conclusions
hold.
\begin{enumerate}
\item Every maximizing sequence of length-one exact problems satisfies
\begin{equation}
\sup_{0\le r\le1}
 \big|f_i(r)-f_i(0)-d_*r\big|\longrightarrow0 .
 \label{a03:linear-profile}
\end{equation}
Only the relative profiles occur; no bound on \(f_i(0)\) is required.
\item For every fixed finite ordered grid of target depths ending at one,
there are nested true systems at actual depths tending to those targets.
Their cumulative horizons tend to \(kr\), and their segment horizons,
in their own length units, tend to \(k\).  Both the coarsened and remaining
problems used at each step are maximizing after their individual length
normalizations.  Full layer assignments, priors and incidence witnesses
can be retained in addition to the final path incidences.
\item In version 2 with \(\ell<\infty\), for every \(\zeta>0\) one may
choose the base maximizing sequence and these paths so that cumulative
narrowness at a target \(r\), and narrowness increments between targets,
equal the corresponding linear values with slope \(\ell\), up to
\(O(\zeta)+o_i(1)\).  The base sequence has terminal optimal narrowness
supremum at most \(\ell+\zeta\).
\item If \(\ell=\infty\), every newly obtained maximizing sequence admits,
on subsequences, true terminal systems with relative horizons tending
to \(k\) and relative narrowness tending to infinity.  Thus arbitrarily
large narrowness increments can be required on each of finitely many
successive positive-length segments.
\end{enumerate}
\end{proposition}
\begin{proof}
\par\smallskip\noindent\textbf{Relative profiles.}\quad

We first establish the relative-profile assertion,
\begin{equation}
f_i(r)-f_i(0)\longrightarrow d_* r
\end{equation}
uniformly.

Take any uniform subsequential limit of the relative profiles by their increment bounds. If it has a secant on an interior \([a,b]\) of slope strictly below \(d_*\), fix \(0<d''<d_*\) still larger than the secant slope. Minimize \(f_i(x)-d''x\) on \([a,b]\), obtaining \(r_i\) bounded away from \(a\) for large \(i\) by the strict secant comparison and Lipschitz bound. At the coarsened endpoint \(r_i\) we have the slope cap \(d''\) back to \(a\); there is also a global terminal cap strictly less than one by the ambient Lipschitz bound. Indeed from \(x<a\),
\begin{equation}
f_i(r_i)-f_i(x)\le r_i-x-(1-d'')(r_i-a).
\end{equation}
Let \(G_i\) be this coarsened problem.  We claim that \(H(G_i)=0\).
Otherwise fix this \(i\), and let \(C_i,c_i\) be the constants in
\Cref{lem:terminal-extension} for its global terminal cap.
Choose
\[
 0<e<\min\bigl(H(G_i),c_i r_i,(r_i-a)/C_i\bigr)
\]
and a true terminal system with horizon \(h<H(G_i)+e/4\).
Since \(h\ge H(G_i)>e\), terminal extension applies.  It gives
a known atlas with horizon \(h-e\) at a depth \(t_i\) satisfying
\[
 0<r_i-t_i\le C_i e<r_i-a,
\]
so \(t_i>a\).

Normalize through this atlas to obtain a homogeneous remaining
problem \(F\) of length \(L_F=r_i-t_i\).  The terminal density
equality and intermediate density upper bound in
\Cref{lem:normalization} give, for its profile \(f_F\),
\[
 f_F(L_F)-f_F(s)
 \le f_i(r_i)-f_i(t_i+s)
 \le d''(L_F-s),\qquad 0\le s\le L_F.
\]
The second inequality uses only the slope cap on \([a,r_i]\),
since the entire remaining interval lies there.  The angular cap
also passes through normalization, so \(F\in C(d'')\).
By \Cref{a03:known-composition},
\[
 H(F)\ge H(G_i)-(h-e)>3e/4>0,
\]
contradicting the definition of \(d_*\).  All choices in this
argument are within the fixed exact problem \(G_i\); neither
\(e\) nor its extension constants need be uniform in \(i\).

Using these zero outer costs back in the original problems gives \(H(E_i)\le k(d_i)(1-r_i)\), contradicting maximization since \(k>0\). If \(d_*=0\), a lower secant is already excluded by monotonicity.  Thus in every case no lower interior secant exists. Together with the endpoint cap (in the limit) and continuity this forces the linear profile, proving the assertion.

\par\smallskip\noindent\textbf{Nested true systems.}\quad

Fix a target \(0<r<1\) on a maximizing sequence.  Choose
\(\epsilon_i>0\) tending to zero sufficiently slowly relative to
the uniform profile errors, and minimize
\(f_i(x)-(d_*+\epsilon_i)x\) on \([0,r]\).
The linear approximation forces a minimizer \(r_i\to r\).
At this actual endpoint the terminal profile cap back to the start
is \(d_*+\epsilon_i\).

The coarsened problems at \(r_i\), each normalized to length one,
form a maximizing sequence.  Indeed, their outer costs before
rescaling satisfy
\begin{equation}
H_{\rm out}\le k(d_*+\epsilon_i)r_i,\qquad
 H(E_i)\le H_{\rm out}+k(d_i)(1-r_i),
\end{equation}
so \(H_{\rm out}/r_i\to k\).

Take true systems on these outer problems with horizon exponents tending to \(k\) in relative units.

Normalize the original finer problems through these packets.
The resulting homogeneous remaining problems are again maximizing
when normalized to length one, with base \(N_0^{1-r_i}\): their
costs in old units lie between
\(H(E_i)-(k+o_i(1))r_i\) and \(k(d_i)(1-r_i)\).
These conclusions persist on further subsequences with the same
choices.

We can therefore iterate over any fixed finite ordered grid of
target depths ending at one.  At each step use the next target
fraction of the remaining length, choosing that fixed fraction
from the limiting remaining length.  The maximizing property just
proved permits the next application; at fraction one use the full
remaining maximizing sequence.  This gives nested true systems
whose actual depths tend to the targets and whose segment horizons
have the required limits, with the final node at the true endpoint.

Each normalization divides by its individual length, and returning
to old units multiplies by that length, not merely its limit.
Only finitely many subsequences of \(i\), and of the individual
exact sequences, are needed for a fixed grid.

Intermediate true packets lift through true parents by the equality
case of \Cref{lem:normalization}.  At each stage homogenize and
prepare them with the bounds following \Cref{def:true-chart}.
Keep their full layer assignments, priors and incidence witnesses,
as well as the final path incidences common to all layer selections.
The final path alone need not saturate an earlier chart.
In a parent, for example, the trajectory prior remains conditioned
on its full assignment after further incidence selections.
Child trajectory sets use this fixed prior and lift to subsets
of the parent assignment.

Labels may include their whole histories; list knowledge composes via the parent's knowledge and the bounded overlaps of the corresponding hidden grids once chart and exact base are specified. Time lengths and spatial Jacobians multiply. All these paths for a fixed finite grid and \(i\) use actual exact subsequences and their subpower slacks.

\par\smallskip\noindent\textbf{Narrowness optimization.}\quad

In version 2 with \(\ell<\infty\), fix a small tolerance \(\zeta>0\)
and choose a whole base maximizing sequence with
\(\omega\le\ell+\zeta\).  At each step of the fixed finite-grid
construction, including the final segment, choose a true system whose
relative horizon exponent tends to \(k\) and whose relative narrowness
has \(\liminf\) at least \(\ell-\zeta\).  To make this choice, apply
the definition of \(\ell\) to the newly coarsened maximizing sequence
in its length-one units: choose an attainable \(\limsup\) large enough,
then pass to a subsequence.  The new remaining problems continue to
maximize by the preceding composition bounds.

Compose these systems through the final depth.  True lifting gives a
terminal true system with cumulative horizon exponent tending to
\(k\), so its cumulative narrowness has \(\limsup\le\omega\) for the
chosen base sequence.  Narrowness exponents add, with each relative
exponent multiplied by its \emph{actual} segment length in original
units.  The segment lower bounds therefore give the cumulative lower
bound at a node of target \(r\); subtracting the lower bounds for the
subsequent segments from the terminal upper bound
\(\ell+\zeta+o_i(1)\) gives the upper bound.  Thus the cumulative
narrowness is \(\ell r+O(\zeta)+o_i(1)\), and the corresponding
increments have the same linear behavior.  Cumulative horizon
exponents tend to \(kr\).  No single exact limiting problem is
assumed to realize these objectives.

If instead \(\ell=\infty\), every newly obtained maximizing sequence
admits, on a subsequence, true terminal systems with relative horizons
tending to \(k\) and relative narrowness tending to infinity.
Choose arbitrarily high narrowness thresholds and diagonalize over
\(i\) in the definition of \(\omega\).

\end{proof}

\begin{corollary}[Depth-zero normalization and vanishing-depth nodes]
\label{a03:zero-depth}
A sequence of depth-zero, horizon-zero known normalizations of a
maximizing sequence is again maximizing, provided the normalizations
have the actual subpower success and terminal knowledge required in
\Cref{lem:normalization}.  This assertion does not require bounded
absolute density exponents.

On any maximizing sequence one may choose true nodes at positive depths
\(b_i\to0\) whose horizon exponents tend to zero, retaining the
relative-profile estimates at those nodes in the slowly diagonalized
sense.
\end{corollary}
\begin{proof}
Let \(F_i\) be a depth-zero normalized problem of length one.
Normalization preserves \(C(d_i)\), and lifting gives
\begin{equation}
H(E_i)\le H(F_i)\le k(d_i).
\end{equation}
Both outer bounds tend to \(k\).  Thus \(F_i\) is maximizing, and
\Cref{a03:linear-profile} applies to \(f_{F_i}(r)-f_{F_i}(0)\), without
extracting a limit of either absolute term.

For every fixed \(b\in(0,1)\), the construction in
\Cref{prop:critical-path} gives actual depths \(b_i\to b\), relative
profile errors tending to zero, and horizons \(kb+o_i(1)\).
Apply this statement successively with \(b=1/m\).  For each fixed \(m\)
first take the outer index far enough and retain the required exact
subsequences.  Then increase \(m\) sufficiently slowly that the depth
error and horizon error are at most \(1/m\), and that the profile errors
are as small relative to \(b\) as subsequently required.  Each choice
still uses actual exact problems and their true systems.  It gives
\(b_i>0\), \(b_i\to0\), and horizon tending to zero.  Additional finite
requirements, such as the narrowness choices in
\Cref{prop:critical-path}, can be included before each increase of \(m\).
\end{proof}

\begin{lemma}[Excluded improvements]\label{a03:excluded-improvements}
There are two forbidden outputs of known atlases at an interior path node (including the parent where a short piece starts) as \(i\to\infty\) on the chosen sequence and fixed grid. Exponents refer cumulatively to the original problem; one can homogenize within exact problems. Write \(r_i,h_i\) for that node's depth and horizon exponent. Knowledge is required by the original \(p_L\).  Path incidences alone
may be used, with subpower coverage and slacks in each exact problem.
The gains \(\Delta_i>0\) have a uniform positive lower bound along
the outer sequence.
\begin{enumerate}
\item A thickness-depth improvement by \(\Delta_i>0\) beyond \(r_i\), still leaving positive length, with cumulative time exponent at most \(h_i+k\Delta_i/2\).
\item In the finite-\(\ell\) width optimization, extra narrowness by \(\Delta_i\gg\zeta>0\) over the node's, keeping exactly the node depth and horizon.

\end{enumerate}
\end{lemma}
\begin{proof}
The first on an infinite subsequence would give by normalization and lifting
\begin{equation}
H(E_i)\le h_i+k\Delta_i/2+k(d_i)(1-r_i-\Delta_i)\le k-k\Delta_i/2+o_i(1),
\end{equation}
impossible. In the second the remaining problems after the output atlas still maximize, by the same composition reasoning (they have remaining cost at least \(H(E_i)-h_i\) before length normalization). Take their true terminal packets with relative narrowness \(\ge\ell-\zeta\) in liminf and optimal limiting horizon along a further subsequence. Lifting now gives terminal horizon tending to \(k\) and cumulative narrowness beating \(\ell+\zeta\), by the previously built segment lower bounds and the extra gain (e.g. with lower bound \(>4\zeta\)). This violates the base sequence's upper bound. This explains why mere approximate saturation or approximate label knowledge is not enough.

\end{proof}

\subsection{A tangent diagonal with an original-resolution exclusion}

We next extract short portions of these paths. The local estimates
will use tangent units, while their final outputs must still violate
an exclusion in the original exact problem.

Take stages \(n\to\infty\) and a tangent depth unit \(\epsilon_n\to0\). For fixed \(n\) put disjoint successive windows of length \(D_n\epsilon_n\) in an interior interval bounded away from depths zero and one, using \(\gtrsim1/(D_n\epsilon_n)\) windows. Their starts are potential parent nodes. Subdivide each with a finite mesh of spacing \(o_n(1)\) in relative units \(s=(\text{offset from start})/\epsilon_n\); form true paths for this whole grid, including the terminal depth.

Use a finest indicated coarse velocity precision of depth \(\epsilon_n U_n\) in the \textbf{original} problem coordinates, where \(U_n,D_n\to\infty,\ U_nD_n\epsilon_n\to0\).

Choose a forbidden gain threshold \(\gamma_n\epsilon_n,\ \gamma_n\to0\) positive, and path optimization tolerance \(\zeta\ll\gamma_n\epsilon_n\) as above. Take \(i\) sufficiently large in the resulting sequence, then take a sufficiently late sample \(N_0\) of its exact path subsequence and use
\begin{equation}
N=N_0^{\epsilon_n}\to\infty,\qquad K=N^{o(1)}
\end{equation}
along the stages. All node depths, horizon exponents and profiles needed on the finite grids can use errors \(o_n(\epsilon_n)\) in original exponent units by the preceding limiting results and this order; cumulative narrowness in version 2 also uses the \(O(\zeta)\) tolerance.

Thus later labels at specified \(s\) mean rounded path nodes with error tending to zero in tangent exponent units. Additional analytical grids in fixed finite exponent ranges can increase sufficiently slowly. \(K\) denotes changeable subpower factors \emph{on the diagonal}, possibly much too large as slacks for an output.


At a fixed stage \(n\), choose the outer index \(i\) large enough that
\Cref{a03:excluded-improvements} excludes the buffered gain
\(\gamma_n\epsilon_n\) for this finite grid.  Fix its exact path
subsequence.  Let
\begin{equation}
K_{\mathrm{path}}=N_0^{o_{N_0\mid i}(1)}
\end{equation}
include its true-path geometric and list bounds, cumulative products,
preparation costs and reciprocal selected path mass.  Retain an
original tempered final path submeasure of mass
\(\mathfrak p_{\mathrm{path}}\ge K_{\mathrm{path}}^{-1}\) as a reference,
together with all earlier layer assignments and witnesses.

\begin{lemma}[The finite-sample exclusion]\label{a03:finite-guard}
Fix this \(n,i\), exact path subsequence, and any finite original-coordinate
complexity exponent \(M^\#\).  For all sufficiently late \(N_0\) there
is no output known atlas satisfying all of the following:
\begin{enumerate}
\item One of the improvements in \Cref{a03:excluded-improvements},
with its buffered gain at least \(\gamma_n\epsilon_n\), with thickness
depth bounded away from the terminal endpoint.  A width improvement
keeps the actual node depth and time length and holds uniformly on its
output labels.
\item Counts, coefficient magnitudes, reciprocal widths and lengths,
piece and condition counts are at most \(N_0^{M^\#}\), with degree per
condition at most \(M^\#\).
\item All geometric slacks and original \(p_L\)-list sizes are at most
\((K_{\mathrm{path}}\log N_0)^n\), and the original incidence mass is
at least \(\mathfrak p_{\mathrm{path}}/\log N_0\).
\end{enumerate}
\end{lemma}
\begin{proof}
Otherwise there are infinitely many such outputs along this one exact
sequence.  The fixed complexity bound makes them uniformly tempered.
Their slacks, list sizes and reciprocal masses are subpower in the
original resolution \(N_0\), since \(n,i\) are fixed.  Homogenize the
output geometry in logarithmic order and extract limiting horizon and
thickness-depth exponents.  The strictly buffered improvement persists.
For a width improvement its uniformity on labels, relative to the
homogeneous old widths, ensures that it persists as well.  Replacing
a thickness exponent by its limit changes thickness by an
\(N_0^{o(1)}\) factor; both the value and curvature requirements are
preserved after enlarging the actual subpower slack.  These outputs
therefore form an actual known atlas on an exact subsequence, excluded
by \Cref{a03:excluded-improvements}.
\end{proof}

Let \(M_{\mathrm{path}}\ge1\) bound all original and path complexity
exponents, including law complexity, degrees and coordinate changes.
It is fixed within this exact problem before the sample is chosen.
Declare, for example,
\begin{equation}
M^\#=2^{(M_{\mathrm{path}}+n)^2}.
 \label{a03:complexity-guard}
\end{equation}
The order of the choices is
\begin{equation}
\text{finite stage and grid}
 \ \longrightarrow\
 \text{outer exact problem and its path}
 \ \longrightarrow\
 (M_{\mathrm{path}},M^\#)
 \ \longrightarrow\
 N_0 .
 \label{a03:limit-order}
\end{equation}
Choose \(N_0\) beyond the cutoff of \Cref{a03:finite-guard} for this
already declared \(M^\#\).  It may be taken arbitrarily late there.
If \(G_n\) denotes the total number of relevant grid nodes, impose also
\begin{align}
 \frac{G_n}{\log N_0}&\le\frac1n,
 &
 N_0^{-1}K_{\mathrm{path}}\log N_0&\le\frac1n,
 \label{a03:mass-diagonal}\\
 \frac{\log\big((K_{\mathrm{path}}\log N_0)^n\big)}
      {\epsilon_n\log N_0}&\le\frac1n .
 \label{a03:slack-diagonal}
\end{align}
Each condition is possible in the fixed exact sequence.  All other
finite-stage approximation requirements are enforced at the same time.
In particular the required original profile and horizon errors are
\(o_n(\epsilon_n)\), and the named path factors and logarithms are
\(K=N^{o(1)}\) on the resulting tangent diagonal.  Analytical grids in
fixed finite exponent ranges can be increased sufficiently slowly.

There will be a countable menu of fixed candidate scenarios.  A
scenario will be identified only after the diagonal and a successful
subsequence have been chosen.  The implementation below proves that,
for every fixed scenario \(\mathcal S\), the original-coordinate
complexity of every possible output is bounded by a polynomial
\(P_{\mathcal S}(M_{\mathrm{path}})\), uniformly over the selected tests,
masks and posterior rules.  This fits the predeclared guard eventually.
Indeed, if \(P_{\mathcal S}(M)\le C_{\mathcal S}(1+M)^{d_{\mathcal S}}\),
then
\begin{equation}
\sup_{M\ge1}
 \{d_{\mathcal S}\log_2(1+M)-M^2\}<\infty
\end{equation}
implies \(P_{\mathcal S}(M)\le2^{(M+n)^2}\) for all \(M\ge1\) once
\(n\) is large enough depending only on \(\mathcal S\).  Likewise an
actual slack or list power fixed by \(\mathcal S\) is eventually less
than \(n\).  We discard finitely many stages of the successful
subsequence; we do not alter the guard or resample those stages.
There is no claim that \(H\) is preserved merely by a tangent limit.

\subsection{Preparing the tangent laws}

The path geometry gives the scales of the local problems.  We now prepare
the measures on which the local arguments will find improvements.  Two
operations have different roles.  Density and angular exceptions are
removed simultaneously at every potential parent before we choose an
entropy window.  After that choice the parent probabilities and velocity
palettes are fixed; graph preparations and later selections restrict those
laws without renormalizing them.

\begin{proposition}[Prepared tangent data]\label{prop:tangent-palette}
Assume \(\ell<\infty\), and use the guarded diagonal constructed above.
One can retain all but \(o(1)\) of the relative mass of the original final
path and choose an interior parent window with the following properties.
Write \(d=d_*\), \(N=N_0^{\epsilon_n}\), and use \(K=N^{o(1)}\) for
analytical estimates.
There is one prepared path probability
\(\mu_{\mathrm{prep}}=\sum_Ap_A\lambda_A^0\), where
\(p_A=\mu_{\mathrm{prep}}\{A\}\) and \(\lambda_A^0\) is its
conditional probability in parent \(A\).  On the retained parents,
\(\lambda_A^0\le K\nu_A\times dt\), where \(\nu_A\) is the prior
conditioned on the full parent trajectory assignment and \(dt\) is
uniform parent time.  Subsequent omissions restrict this mixture and
leave its reference weights and probabilities fixed.
\begin{enumerate}
\item At every required bounded tangent depth \(s\), the true label
\(Q_s\) has thickness \(a_s=N^{-s}\), time length
\(q_s\sim_{\log,N}N^{-ks}\), and spatial determinant
\(J_s\sim_{\log,N}q_s^jN^{-\ell s}\) in version~2; omit the last
factor in version~1.  Its largest width is comparable to \(q_s\) in
logarithmic order, and its full assigned trajectory mass satisfies
\(\nu_Q\sim_{\log,N}a_s^dJ_s\).
\item For the restricted incidence laws, pure hidden and native-signal
bins of width \(p\) have the instantaneous joint base-density ceiling \(Kp^d\) on the required
finite grids.  In unit coordinates of a \(Q_s\)-box, the base marginal
has density at most \(Kq_s\nu_Q\), with masses still measured in
parent time.
\item Let \(v=V\) in version~1 and in the case \(\ell=0\).  For
\(0<\ell<\infty\), let \(Z,Y\) be the major and minor parent
coordinates and put \(v=Z'\).  The fixed conditional parent palette
\(\pi_A\), the velocity marginal of \(\lambda_A^0\), satisfies
\(\pi_A(B_p)\le Kp^S\).  If \(E\) denotes the end history,
then before likelihood deletion the required velocity bin satisfies
\[
 \sum_Ap_A I_{\lambda_A^0}([v];E)=o(\log N).
\]
Thus the information bound is averaged with the fixed parent weights.
\item Required finite-mesh point states and path labels have lists
at the original \(p_L\) bounded by fixed powers of actual path factors
and logarithms.  Sufficiently deep histories leave only \(K\) options
for these point states.  Each fitted test has at most \(K\) local graph choices approximating
\(w\) to \(Ka_s\) on retained incidences in their cores.  After composition with the native test, their
fixed-order derivatives are bounded by \(K\) on the used normalized
interiors.  These graph bounds and capped full-layer incidence witnesses
can be prepared before arbitrary later inverse-subpower selections.
\end{enumerate}
When \(0<\ell<\infty\), major-coordinate fullness and the coefficients
of the tilted-box rows can be prepared with actual \(N_0\)-subpower
bounds.  The joint and conditional consequences of these fixed laws are
stated in \Cref{a03:finite-state-conditioning} below.
\end{proposition}

\begin{proof}
\leavevmode\par
\paragraph{Coordinates and masks at every potential parent.}
For now retain all potential parent nodes in the finite grid.  If \(A\)
is a label of cumulative depth \(b\), let \(q_A\) be its original time
length, \(J_A\) its spatial determinant, and
\(\nu_{\mathrm{old}}(A)\) its full assigned trajectory mass in its
old world.  Saturation gives
\[
 \nu_{\mathrm{old}}(A)
 \sim_{\log} n_{\mathrm{old}}(b)J_A.
\]
Normalize its base coordinates as in \Cref{lem:normalization}.  The
parent prior \(\nu_A\) is the old trajectory prior conditioned on
this full assignment, and parent time is uniform on \([0,1]\).
In version~2 rescale the hidden coordinate so that its parent derivative
unit is \(B=1\).  Its value need not be bounded in tangent units; the
native signal \(X=P_*\) has the native bounds.

At width \(p=N^{-r}\), a parent hidden bin corresponds, at fixed base
and label, to boundedly many old bins at depth \(b+\epsilon_n r\).
Remove the original high-density exceptions at these depths, for all
required parents and finite grids.  The within-world density factor
under parent normalization is \(J_A/\nu_{\mathrm{old}}(A)\).
Consequently the remaining unnormalized density is at most
\[
 K\,\frac{J_A}{\nu_{\mathrm{old}}(A)}
       n_{\mathrm{old}}(b+\epsilon_n r)
 \le K\,\frac{n_{\mathrm{old}}(b+\epsilon_n r)}
                 {n_{\mathrm{old}}(b)}
 \le Kp^d.
\]
The last inequality uses the relative-profile errors
\(o_n(\epsilon_n)\) fixed when choosing \(i\).  The native derivative
and curvature bounds give the same ceiling for \(X\)-bins.  These
are ceilings for restricted measures, so they remain valid under all
subsequent deletions.

Prepare angular ceilings on the same finite collection of parents.
Refinement of the original terminal observation by a known parent label
has typical density at least \(n(L)/K_{\mathrm{path}}\) under the
unnormalized final path measure.  Remove the exceptional denominators
and the old joint angular exceptions before projecting velocity.  A
normalized velocity cell of side \(p\) pulls back to an old velocity
ball of radius \(Kp\).  In the positive-narrowness case the omitted
minor component is already much smaller than this radius, since the
parent is at an interior original depth.  For every tested cell the
remaining angular numerator is therefore bounded by \(Kp^S\) times
the original refined terminal denominator.  Integrating gives the
corresponding bound relative to the parent masses before these masks.
No exceptional portion is included in this numerator sum.

We also prepare the spatial fullness that will be needed in original
units.  For \(0<\ell<\infty\), use the cumulative principal axes of
each potential parent.  Its minor spatial scale per unit time is
smaller than the major scale by an actual power of \(N_0\).  By
\Cref{a01:largest-axis}, cumulative child largest widths per unit
time are full up to actual subpower factors.  Make these prunings
simultaneously for the finitely many required ancestor--child pairs.
For a descendant at a specified relative path depth, let \(D_s\) be
its actual relative spatial matrix and \(q_s\) its actual relative
time length.  Then
\[
 \rho_s=\|D_s^{\mathsf T}e_Z\|
       \ge q_s/K_{\mathrm{path}}.
\]
Indeed \(\|D_s\|/q_s\) is actually subpower bounded.  If its major
row were deficient by a fixed power, composition with the parent's
matrix, whose minor scale already has a fixed-power deficiency, would
make the cumulative child largest scale deficient as well.  This
contradicts the fullness just prepared.

All density exceptions have power-small mass at fixed tolerances in
each exact problem.  The original order \(n,i,N_0\) allows the finite
grids, tolerances, and fullness prunings to be chosen with total loss
\(o(1)\) relative to the reference path.  At this stage no entropy
window has been selected.

\paragraph{The fixed probability law and the entropy window.}
Normalize the surviving path once and call its probability law
\(\mu_{\mathrm{prep}}\).  For a potential parent \(A\), let \(s_A\)
be its surviving success fraction in \(\nu_A\times dt\).  For
\(s_A>0\), set
\[
 \lambda_A^0
   =\frac{(\nu_A\times dt)|_{\mathrm{prepared\ path\ in}\ A}}{s_A},
 \qquad p_A=\mu_{\mathrm{prep}}\{A\}.
\]
At any one parent node,
\(\mu_{\mathrm{prep}}=\sum_Ap_A\lambda_A^0\).
The weight \(p_A\) is proportional to its old world weight times
\(q_A\nu_{\mathrm{old}}(A)s_A\).

Parents with too small a success fraction, or too small a surviving
denominator relative to their pre-mask mass, have negligible total
weight: sum the former costs with
\(\sum_A\omega_{\mathrm{old}}q_A\nu_{\mathrm{old}}(A)\le1\),
and the latter with the pre-mask parent masses.  These parents may be
omitted when using the chosen window.  On the parents retained for
estimates,
\begin{equation}
 \lambda_A^0\le K\nu_A\times dt,
 \qquad \pi_A(B_p)\le Kp^S,
 \label{a03:prepared-parent-law}
\end{equation}
where \(\pi_A\) is the velocity marginal of \(\lambda_A^0\).
The division by \(s_A\) costs only \(K\), so the pure ceilings above
also hold for these conditional probabilities.

Choose the window using the single law \(\mu_{\mathrm{prep}}\),
before any further graph or likelihood deletion.  Let \(V_*\) be
the indicated original velocity bin, of depth \(\epsilon_n U_n\).
For disjoint successive windows, the increments
\[
 I(V_*;\hbox{history through the window end}
           \mid\hbox{history through its start})
\]
are nonnegative and sum to at most \(H(V_*)\), by the chain rule
for nested histories.  The original velocity range gives
\(H(V_*)\le O(U_n\log N)+o(\log N)\).  There are
\(\gtrsim(D_n\epsilon_n)^{-1}\) windows, so one has increment
\(o(\log N)\), since \(U_nD_n\epsilon_n\to0\).

Use the normalized coordinates of this window's parent.  In version~1
and for \(\ell=0\), put \(v=V\); for positive finite narrowness,
put \(v=Z'\).  Given the parent, the required \(v\)-bin has at most
\(K\) possibilities from \(V_*\), and hence its conditional mutual
information with the end history is larger by at most
\(\log K=o(\log N)\).  For \(\ell=0\), both parent scales per
unit time are full to tangent factors \(K\), which suffices for this
ambiguity estimate; positive narrowness uses the major-coordinate
fullness prepared above.  The requested velocity grid is no finer
than \(N^{-U_n}\) in its own units.

Fix \(p_A\), \(\lambda_A^0\), and \(\pi_A\) from now on.  Omitting
the exceptional parents and imposing graph masks restricts this
mixture without redefining its reference laws or its entropy estimate.
Substantial coverage will mean positive probability relative to the
original reference path.

\paragraph{Path scales and the two time conventions.}
For a descendant \(Q=Q_s\), put
\(a=N^{-s}\), \(h=N^{-ks}\), and, in version~2,
\(g=N^{-\ell s}\).  The relative-profile and path conclusions give
\begin{equation}
 a_s=N^{-s},\qquad q_s\sim_{\log,N}h,
 \qquad J_s\sim_{\log,N}h^jg,
 \qquad \nu_Q\sim_{\log,N}a^dJ_s.
 \label{a03:tangent-scales}
\end{equation}
In version~1 omit \(g\).  Write \(P_s\) for the fitted true test.
Its actual node thickness \(a_Q\) is comparable to \(a\) in
logarithmic order.  The true value,
derivative, and curvature bounds continue to hold in the actual
units of that node.  The full trajectory assignments remain disjoint
within each time block and nested in the fixed parent prior.

At an exact base, the true value and incidence-derivative bounds
permit only \(K\) hidden bins of width \(a_Q\).  The map to unit
base coordinates \(u\) of \(Q_s\) has Jacobian \(q_sJ_s\);
denote this map by \(\operatorname{base}_Q\).
The pure cap therefore gives
\begin{equation}
 d(\operatorname{base}_Q)_*(\lambda_A^0|_Q)(u)
   \le Ka^dq_sJ_s\,du
   \le Kq_s\nu_Q\,du.
 \label{a03:packet-base-cap}
\end{equation}
This is a cap in parent-time mass.  Dividing by the block length
\(q_s\) gives the cap \(K\nu_Q\) in normalized \(Q_s\)-time.
If the trajectory prior is normalized on the full assignment of \(Q\),
the pure ceiling is \(Kp^d/\nu_Q\) in parent base coordinates.
After changing to unit \(Q\)-base coordinates and unit block time,
it is \(Kp^dJ_s/\nu_Q\).  Here \(p\) remains the hidden-bin width
in parent units; the denominator is the full assignment mass, not the
current selected incidence mass.
The budgets are
\[
 \sum_{Q\ \mathrm{in\ one\ block}}\nu_Q\le1,
 \qquad \sum_{Q\ \mathrm{at\ one\ node}}q_s\nu_Q\le1.
\]
They allow light labels to be removed whenever current incidence
floors are needed, with thresholds chosen below the current selected
mass.  Spatial enlargements by subpower factors only change \(K\).

For positive finite narrowness, subtract a multiple of the \(Z\)-row
of \(D_s\) from its \(Y\)-row to make the rows orthogonal.  The
subtraction coefficient is at most
\(\|D_s^{\mathsf T}e_Y\|/\rho_s\), and the new transverse row has
length \(|\det D_s|/\rho_s\).  Thus the child is a tilted box of
widths \(h,hg\), and its assigned velocities satisfy
\[
 Y'=A_s\cdot(1,v)+O(Kg).
\]
The row \(A_s\) is determined by \(Q_s\) and has actual subpower
coefficients.  Its affine center velocity is controlled by the
whole-block position bounds on occupied trajectories.  In actual
units the transverse width per time is comparable, up to actual
subpower factors, to \(J_s/q_s^2\).

\paragraph{Stable graphs and point-state lists.}

Later arguments need velocity bounds after recording finite-mesh
point states.  We first construct the finite list of possible states;
its size will be the input to the appended-entry floor below.
Use normalized base coordinates \(u\) in \(Q_s\), keeping the hidden
variable \(w\) in parent units with \(B=1\).  Let \(K_g=N^{o(1)}\)
contain the present test bounds.  In version 2 they give
\[
 |P_s|\le K_ga_Q,\qquad
 K_g^{-1}\le |P_{s,w}|\le K_g,\qquad
 |P_{s,ww}|\le K_g/a_Q
\]
on the counted incidences.  First omit labels whose current mass is
too small relative to \(q_s\nu_Q\).  The assignment budget controls
the total loss, and \Cref{a03:packet-base-cap} then gives a base
projection of inverse-subpower volume in every used label.

The polynomial discriminant
\[
 \mathcal D(u)=P_{s,w}^2-2P_{s,ww}P_s
\]
is independent of \(w\).  Its values on this projection are bounded
by \(K\); polynomial Remez therefore bounds \(\mathcal D\) and
its first derivatives by \(K\) on the needed complex enlargements
of the normalized packet box.  Choose a subpower \(H\) with
\(H/K_g^3\to\infty\), and split the tests at
\(\lvert P_{s,ww}\rvert=(a_QH)^{-1}\).
Above this cutoff the affine quadratic center satisfies
\[
 |w-F_{\mathrm{ctr}}(u)|
   =\frac{|P_{s,w}(u,w)|}{|P_{s,ww}|}
   \le K_gHa_Q\le Ka.
\]
Below the cutoff,
\begin{equation}
 2|P_{s,ww}P_s|\le \frac{2K_g}{H}
       =o(K_g^{-2})=o(|P_{s,w}|^2),
 \qquad \mathcal D\ge\tfrac12K_g^{-2}
 \label{a03:palette-discriminant}
\end{equation}
at those incidences.  Partition the normalized box into core
subboxes of inverse-subpower radius, each with a fixed-factor
complex enlargement and a further margin.  Choose the radius so
that the variation of \(\mathcal D\) on each occupied enlargement
is much smaller than this lower bound.  For a quadratic test, each
such enlargement has two stable simple holomorphic branches, real
on its real core.
The branch with the incidence derivative sign approximates \(w\)
to \(Ka\), by \Cref{a02:nearest-quadratic-root}.  For a linear
equation the same argument keeps its coefficient nonzero.

There are \(K\) core/branch entries per label, including affine
center entries.  Delete light entries using the same
\(q_s\nu_Q\) budgets and a smaller reciprocal-subpower threshold.
The cap \Cref{a03:packet-base-cap} supplies inverse-subpower real
base volume for each retained entry.  Compose its graph with the
native test \(P_*\).  The native derivative and curvature bounds
give an error \(Ka\) in \(X\) at the incidences, and \(|X|\le K\)
there.  The composition is an algebraic holomorphic sheet;
\Cref{lem:algebraic-norm} bounds it and its required derivatives
by \(K\) on the used interiors.  In version 1 use the explicit
graph directly.  These bounds now belong to the fixed graphs and
persist under later incidence restrictions.

An end history refining \(Q_s\) localizes its normalized base
coordinates to uncertainty at most \(Kq_{\rm end}/q_s\).
Indeed, nested relative frames have largest width bounded by
relative time length times path slacks, and occupied chart centers
move at bounded speed up to those slacks.  For a prescribed mesh
in \((t,y,X)\), first choose \(a\) much finer than that mesh and
then choose the end depth so that the graph derivative bound times
\(Kq_{\rm end}/q_s\) is also much finer.  Each graph then gives
only \(K\) possible mesh states.  The finite core/branch list
proves the required ambiguity bound before any palette floor is
used.  Cores lie inside the derivative-controlled interiors;
bounded overlapping covers handle their boundaries.

These preparations use arbitrary fixed mesh exponents first,
with sufficient window depth, and then slowly increasing grids.
They may be made after the entropy-window choice with negligible
total loss, without redefining \(\lambda_A^0\), and before arbitrary
later sparse selections.

\paragraph{Original observation lists and full-layer witnesses.}
Finite-mesh point states in \((t,y,X)\), together with the
path labels in the window, have actual-list accuracy given the
original \(p_L\).  Condition on the actual path charts in its
lists.  The base is exact, and the much finer original terminal
hidden bin, with the native derivative and curvature bounds at
the incidence, controls \(X\).  This is actual knowledge in the
original problem; the graph lists separately supply the
\(K\)-ambiguity needed in tangent estimates.

For stored full-layer incidence witnesses, impose the same original
pure-cap masks on the fixed density grids needed in the coordinate
transfers.  Their costs are power-small at every fixed tolerance in
the exact problem.  If a layer label loses at least half its incidence
mass, saturation bounds its old block length times its full trajectory
mass by an actual path factor times the lost incidence mass.  Sum with
the old world weights.  The cost of dropping all such labels from the
final path is therefore negligible when \(N_0\) is sufficiently late.
This prepares capped witnesses in each needed layer.  Their earlier
floors refer to those witnesses; they are not floors for the final path,
and they do not provide several-time palette estimates.

All these preparations use the already fixed diagonal.  Their initial
cap masks retain \(1-o(1)\) of the reference path, and their graph and
witness prunings have negligible additional relative loss.  The parent
weights and conditional probabilities remain the reference laws fixed
at the entropy-window step.
\end{proof}

\subsection{Conditioning finite states and taking time slices}

The prepared law has little information about the velocity in its end
history.  The next lemma converts that information bound to a pointwise
joint ceiling.  It also records how that ceiling passes to a point state
and how much mass must be deleted before it becomes a conditional
ceiling under a later restriction.

\begin{lemma}[Likelihood truncation and finite-state conditioning]
\label{a03:finite-state-conditioning}
As \(N\to\infty\), let \(p_A\ge0\), \(\sum_Ap_A=1\), be mixture weights and let
\(\lambda_A^0\) be probability laws carrying two finite labels
\(E,b\).  Write
\[
 R_A^0(E,b)=\lambda_A^0\{E,b\},\qquad
 \lambda_{A,\mathrm{pre}}(E)=\sum_bR_A^0(E,b),\qquad
 \pi_A(b)=\sum_E R_A^0(E,b).
\]
Suppose
\[
 \sum_Ap_A\sum_{E,b}R_A^0(E,b)
 \log\frac{R_A^0(E,b)}
 {\lambda_{A,\mathrm{pre}}(E)\pi_A(b)}=o(\log N).
\]
Then deleting \(o(1)\) mass from the whole mixture gives joint
submeasures with
\begin{equation}
 R_{A,\mathrm{good}}(E,b)
 \le K\lambda_{A,\mathrm{pre}}(E)\pi_A(b),
 \qquad K=N^{o(1)}.
 \label{a03:palette-joint}
\end{equation}

Let \(R_{A,2}\) be any later submeasure, extended by a label \(d'\)
with at most \(D\le K\) possibilities per \(E\), whose \((E,b)\)
marginal is dominated by \(R_{A,\mathrm{good}}\).  For a subpower
\(K_1\), deleting at most \(D/K_1\) mixture mass gives joint ceilings
against the predeletion \((E,d')\)-masses with multiplier \(KK_1\).
These ceilings sum to any coarser state \(C=C(E,d')\), with the
predeletion \(C\)-marginal as reference weight \(w_A(C)\).
For a further restriction \(T_A\) of the retained measure, deleting
states whose current
mass is below \(w_A(C)/K_2\) costs at most \(1/K_2\) mixture mass
and gives, on surviving states of positive mass,
\[
 T_A(b\mid C)\le KK_1K_2\pi_A(b).
\]
The subpowers \(K_1,K_2\) may be chosen so each deletion is negligible
relative to a specified current inverse-subpower mixture mass.
\end{lemma}
\begin{proof}
Put
\(R_{A,\mathrm{prod}}(E,b)=\lambda_{A,\mathrm{pre}}(E)\pi_A(b)\).
The negative part of the log ratio has expectation at most one
in each parent, since
\[
 \sum_{R_A^0<R_{A,\mathrm{prod}}}
 R_A^0\log(R_{A,\mathrm{prod}}/R_A^0)\le1
\]
by \(x\log(1/x)\le1/e\).  Consequently its positive part has
mixture expectation \(o(\log N)\).  Choose
\(\delta_N\downarrow0\) slowly enough that this expectation is
\(o(\delta_N\log N)\).  Delete the sampled pairs where the log
ratio exceeds \(\delta_N\log N\).  Markov's inequality gives
\(o(1)\) loss in the whole mixture, and the retained joint measure
satisfies
\begin{equation}
 R_{A,\mathrm{good}}(E,b)
 \le N^{\delta_N}\lambda_{A,\mathrm{pre}}(E)\pi_A(b)
 \le K\lambda_{A,\mathrm{pre}}(E)\pi_A(b).
 \label{a03:good-joint-numerator}
\end{equation}
The inequality holds in each retained parent; the small-loss
assertion is only for their weighted mixture.  Any further restrictions
of this numerator leave its reference denominator unchanged.

Let \(R_{A,2}\) be a later submeasure of the retained incidence
law, carrying a datum \(d'\) with at most \(D\le K\) possibilities
per end history.  Its marginal on \((E,b)\) is dominated by
\(R_{A,\mathrm{good}}\).  Define its appended-entry masses before
the next omission by
\[
 m_{A,\mathrm{preomit}}(E,d')
    =\sum_bR_{A,2}(E,d',b).
\]
Keep only entries satisfying the floor
\begin{equation}
 m_{A,\mathrm{preomit}}(E,d')
       \ge \lambda_{A,\mathrm{pre}}(E)/K_1.
 \label{a03:appended-entry-floor}
\end{equation}
The omitted mass is at most
\(D K_1^{-1}\sum_E\lambda_{A,\mathrm{pre}}(E)=D/K_1\)
in that parent's probability units.  Choose the subpower \(K_1\)
so this cost is negligible relative to the selected mixture mass.
For the remaining measure \(R_{A,3}\), the numerator and floor give
\begin{align}
 R_{A,3}(E,d',b)
 &\le R_{A,\mathrm{good}}(E,b)
 \le K\pi_A(b)\lambda_{A,\mathrm{pre}}(E)\notag\\
 &\le KK_1\pi_A(b)m_{A,\mathrm{preomit}}(E,d').
 \label{a03:appended-joint-cap}
\end{align}
Since whole entries were omitted, their retained conditional
palette probabilities are bounded by \(KK_1\pi_A(b)\).

For a coarser state \(C=C(E,d')\), sum
\Cref{a03:appended-joint-cap} with the fixed reference weights
\[
 w_A(C)=\sum_{(E,d'):C(E,d')=C}
             m_{A,\mathrm{preomit}}(E,d').
\]
Then \(R_{A,3}(C,b)\le KK_1\pi_A(b)w_A(C)\).
No omitted bad portion is summed against small palette probabilities.

A further sparse restriction \(T_A\le R_{A,3}\) still has this
numerator bound, but its current denominator can shrink.  Before
using a conditional cap under \(T_A\), choose a new subpower
\(K_2\) and discard states with
\[
 T_A(C)<w_A(C)/K_2.
\]
Their total mass is at most \(K_2^{-1}\sum_Cw_A(C)\) in that
parent, and hence at most \(K_2^{-1}\) in the mixture.  Taking
this below the current selected mass gives, on the surviving states,
\(T_A(b\mid C)\le KK_1K_2\pi_A(b)\).
\end{proof}

Apply the lemma with \(E\) the deep end history and \(b=[v]\) under
\(\lambda_A^0\).  The graph preparation has supplied at most \(K\)
point states per sufficiently deep history, so these may be used for
\(d'\).  To treat several rough grids, append one finer dyadic
\((t,y,X)\) mesh and then sum to the required states.  The reference
weights \(w_A(C)\) are marginals of the preceding restricted measure
and retain its pure cell ceilings.  A new sparse conditional law uses
new denominator pruning; an unnormalized numerator ceiling itself
persists under restriction.  All initial likelihood and graph-ambiguity
masks precede unrestricted sparse selections.  This gives no
several-time independence of \(v\).

The resulting bounds are integrated in time.  The following elementary
estimate gives fixed-time bounds for the trajectory prior while keeping
that prior unchanged.

\begin{lemma}[Retained duration and fixed-time prior caps]
\label{a03:time-slice}
Let \(\nu\) be a trajectory probability and
\(d\lambda(l,t)=G(l,t)\,d\nu(l)\,dt\), where \(0\le G\le K\)
on \([0,1]\).  Fix \(K_2\ge1\), partition time into dyadic intervals
\(I\), and set
\[
 L_I=\left\{l:\int_I G(l,\tau)\,d\tau\ge |I|/K_2\right\}.
\]
Deleting the complementary line--interval pairs costs at most
\(1/K_2\) incidence mass.  Fix \(t\in I\).  If trajectory
constraints \(C_t\) and \(C^+_\tau\) satisfy
\(C_t(l)\Rightarrow C^+_\tau(l)\) for every \(\tau\in I\), then
\begin{equation}
 \nu(L_I\cap C_t)
 \le\frac{K_2}{|I|}
       \int_I\int\mathbf1_{C^+_\tau}(l)G(l,\tau)
                         \,d\nu(l)\,d\tau.
 \label{a03:retained-duration-cap}
\end{equation}
For almost every \(t\), the instantaneous restricted incidence mass
on \(C_t\) is at most \(K\nu(L_I\cap C_t)\).
\end{lemma}
\begin{proof}
On a deleted line--interval pair the incidence duration is less than
\(|I|/K_2\).  Integrating over the probability \(\nu\) and summing
\(\sum_I|I|=1\) proves the deletion bound.  On \(L_I\cap C_t\), the
implication in the hypothesis and the duration floor give
\[
 \int_I\mathbf1_{C^+_\tau}(l)G(l,\tau)\,d\tau
 \ge |I|/K_2.
\]
Integration over these trajectories proves the displayed estimate.
The last assertion follows from \(G\le K\).
\end{proof}

In the prepared tangent problem, choose the time partition much finer
than the finite collection of spatial meshes being used.  Bounded motion
of \((t,y,X)\) and of affine trajectory quantities gives the implication
\(C_t\Rightarrow C^+_\tau\) for fixed enlargements of mesh constraints;
a velocity bin stays fixed on a trajectory.  Thus every integrated pure
or joint cell ceiling for these quantities yields the corresponding
fixed-time active-prior cap by \Cref{a03:retained-duration-cap}.
This step requires no time regularity of \(w\).

After a selection, small conditional successes are omitted using the
fixed mixture weights \(p_A\); integrated exceptional numerators are
removed before further sparse conditioning.  Choose \(K_2\) large
enough that its deletion cost is negligible relative to the current
selected mass.  First treat fixed finite mesh exponents, with time
meshes sufficiently fine, and then increase the ranges slowly.
These operations retain the single diagonal already chosen.  In the
return to actual outputs, substantial coverage continues to mean
positive probability on the original reference path, and the final
pigeonholes may divide by stage constants much smaller than
\(\log N_0\).

\subsection{From local candidates to actual improvements}

The later geometric arguments will find a single improved fit inside
any substantial residual of the reference path.  Such a fit need not
carry enough original mass, and its label need not yet be known from the
original observation.  We first specify the local fit and then prove a
conversion criterion that supplies both properties.
\begin{definition}[Improvement candidates]\label{a03:candidates}
Work in a prepared tangent parent and write \(d=d_*\).  Fix a path
packet \(Q=Q_s\), with actual time length \(q_s\), thickness
\(a=N^{-s}\), and assigned trajectory mass \(\nu_Q\).  The tangent
depth \(s\) remains bounded.  In both definitions below, incidence
mass means residual path mass in the parent, divided by \(q_s\).

A \emph{time-improvement candidate of gain \(c>0\)} is a test of the
proper quadratic model class, at width \(a'=aN^{-c}\), on trajectory
pieces assigned to \(Q\).  Its value and hidden-derivative upper
bounds hold throughout the \(Q\)-block:
\[
 |P|\le Ka',\qquad |P_w|\le K,\qquad |P_{ww}|\le K/a'.
\]
On counted incidences it also satisfies \(|P_w|\ge K^{-1}\), and
these incidences have mass at least
\[
 K^{-1}N^{-dc}\nu_Q.
\]

When \(0<\ell<\infty\), put \(g=N^{-\ell s}\).  A
\emph{strip candidate of gain \(c>0\)} is a fixed row \(R\) such that
\[
 |Y'-R\cdot(1,v)|\le KgN^{-c}
\]
on trajectory pieces carrying incidence mass at least
\(K^{-1}\nu_Q\), in the same convention.  Actual rounded path nodes are
understood in both definitions.

These mass tests may use restrictions of \(\lambda_A^0\) without
renormalization, or the parent trajectory-times-uniform-time prior
restricted to the incidences.  On each used parent the two measures
differ by the fixed factor \(s_A^{-1}\le K\).
\end{definition}

For the whole-block upper estimates, extrapolating a polynomial along
a trajectory from relative duration \(K^{-1}\) costs only a factor
\(K\).  The derivative lower bound is required only on the counted
incidences.

A time candidate improves the hidden-coordinate fit without changing the
packet's time block; a strip candidate improves its transverse velocity
prediction.  The mass in the definition is relative to the packet's full
trajectory assignment.  The next proposition converts candidates in every
substantial residual into one original known atlas.

\begin{proposition}[Conversion of candidates to actual improvements]
\label{prop:candidate-upgrade}
Consider the guarded tangent diagonal of
\Cref{prop:tangent-palette}.  Suppose that, on a further subsequence
of every subsequence carrying any substantial residual of the reference
path, that residual supplies a time-improvement or strip candidate
in the sense of \Cref{a03:candidates}.  Its depth remains in a bounded
tangent range and its gain has a fixed positive finite limit, or is
a fixed positive number.

Then one of the original-resolution outputs excluded by
\Cref{a03:finite-guard} exists.  Its original success is at least
\(\mathfrak p_{\mathrm{path}}/\log N_0\), and its lists and geometric
slacks are bounded by a fixed scenario-dependent power of
\(K_{\mathrm{path}}\log N_0\).  Its complete original-coordinate
complexity, including list tables, is bounded by a polynomial in
\(M_{\mathrm{path}}\) depending only on that fixed scenario.
Consequently the candidate hypothesis is impossible on this diagonal.
\end{proposition}
The proof has four steps.  A finite greedy assignment first obtains
substantial coverage by one fixed candidate scenario.  Strip coverage is
converted directly into narrower spatial charts.  For time coverage,
local graph comparison groups the candidate tests into short lists and
then supplies the improved chart tests.  Finally, the following recovery
lemma transfers those lists from the analytical submeasure to the original
law.  The lemma is stated separately because it is also used in the
unbounded-narrowness argument.

\begin{lemma}[Original-law recovery of finite lists]\label{a03:actual-lists}
Let \(\mu\) be an original unnormalized tempered path measure, and let
\(\nu\le\mu\) be an arbitrary measurable submeasure.  Fix a finite
geometric output-label map \(A\), defined on the original law and
agreeing on \(\nu\) with the labels used in the analytical construction.
A failure label may be used off the eligible pieces and is never
included in a list.  Write the original observation as
\(O=(x,o)\), where \(x=(t,y)\) is exact base and \(o\) contains its
original observed discrete entries, including the hidden-coordinate
bin and world.

Suppose a measurable rule \(L(x,o)\), with
\(|L(x,o)|\le D_{\mathrm{list}}\), satisfies
\begin{equation}
\nu\{A\in L(O)\}\ge m.
 \label{a03:masked-success}
\end{equation}
Retain all original observation tags, path tags and output labels in
the joint pushforward law \(\lambda\) of \(\mu\), without making
unobserved tags available to the list rule.  If \(F_\rho\lambda\) is
uniform base-cell averaging at mesh \(\rho\) and
\(\|\lambda-F_\rho\lambda\|_{\mathrm{TV}}\le e\), there is a list table
\(L_C(o)\), indexed only by the base cell and original observed
entries, with the same list size and
\begin{equation}
\mu\{A\in L_{C(x)}(o)\}\ge m-2e.
 \label{a03:cell-list-success}
\end{equation}
If the density, ranges, copies and geometric predicates of this
original joint law have complexity bounded by a fixed polynomial in
\(M_{\mathrm{path}}\) and fixed construction constants, the mesh for
\(e=O(N_0^{-1})\) and the resulting table have complexity bounded by
another such polynomial.  Neither bound depends on the analytical
complexity of \(\nu\) or of \(L\).
\end{lemma}
\begin{proof}
For the fixed geometric map \(A\), submeasure domination gives
\begin{equation}
\mu\{A\in L(O)\}\ge\nu\{A\in L(O)\}\ge m.
 \label{a03:original-dominance}
\end{equation}
Use total variation in the convention of supremum over measurable
events.  In each base cell \(C\) and observed discrete entry \(o\),
choose the \(D_{\mathrm{list}}\) output labels with greatest total
\(\lambda(C,o,\cdot)\)-mass; ties are broken in a fixed ordering.
On \(F_\rho\lambda\) this label-weight vector is independent of the
position in \(C\).  Thus, writing \(E_L\) for the list success event,
\begin{align*}
 \lambda(E_{L_C})
 &\ge (F_\rho\lambda)(E_{L_C})-e\\
 &\ge (F_\rho\lambda)(E_L)-e\\
 &\ge \lambda(E_L)-2e
 \ge m-2e .
\end{align*}
All auxiliary tags remain joint data in this calculation.  The
maximization conditions only on \(C,o\), not on the unknown output
label.  The original hidden-bin boundaries need not align with \(C\):
the original bin is kept as a discrete tag during averaging and its
actual value is used when the table is evaluated on the original law.

Here is the quantitative total-variation bound, including the effect
of marginalizing hidden trajectory coefficients.  In original
coefficient coordinates replace the spatial intercept by \(y(t)\).
The resulting variables are \((x,\xi)\), with
\(x=(t,y(t))\) and \(\xi\) consisting of \(V\) and the remaining
coefficients; this change has Jacobian one.  Include the finite
internal copies.  Before integrating out \(\xi\), regard the density
as a vector indexed by all retained discrete tags, extending it by
zero outside its support.

The original tempered density is piecewise constant on its defining
semialgebraic pieces.  Fix the other variables and translate one coordinate of \(x\) by \(h\).
A nonzero polynomial predicate of degree \(D\) has at most \(D\) roots
on that axis; an identically zero restriction has no sign changes.
The weighted density and its tags therefore have at most the sum of
these degrees many possible changes.  Boolean priority rules and
finite table assignments add no new boundary locations.  If \(H\)
bounds the density, \(T\) the number of predicates, \(D\) their degrees,
\(R\) the coordinate ranges and \(J\) the copy count, integration over
the other variables gives
\begin{equation}
\|\tau_h f-f\|_{L^1}
 \le C H(TD+1)J(1+R)^q |h|,
 \label{a03:axis-variation}
\end{equation}
where \(q\) is fixed by the ambient coefficient dimension.  Marginalizing
\(\xi\) contracts \(L^1\).  Averaging over pairs of points in a base cell,
and changing one coordinate at a time, now gives
\begin{equation}
\|\lambda-F_\rho\lambda\|_{\mathrm{TV}}
 \le N_0^{B_{\mathcal S}(M_{\mathrm{path}})}\rho ,
 \label{a03:joint-tv}
\end{equation}
after absorbing fixed constants.  Here \(B_{\mathcal S}\) is a
polynomial in the old complexity bound for every fixed construction
scenario \(\mathcal S\).  This argument applies to the original
geometric predicates and path law; no analytical mask is inserted.

Choose \(\rho=N_0^{-D_{\mathrm{mesh}}}\) with
\(D_{\mathrm{mesh}}\ge B_{\mathcal S}(M_{\mathrm{path}})+2\).
The error is then \(O(N_0^{-1})\), with room to spare.  If the original
base ranges are bounded by \(N_0^{P_0(M_{\mathrm{path}})}\), the number
of base cells is at most
\begin{equation}
N_0^{(j+1)(D_{\mathrm{mesh}}+P_0(M_{\mathrm{path}}))+O(1)}.
\end{equation}
The number of observed bin entries, worlds and possible labels is
also polynomial.  Multiplying these counts gives the asserted
polynomial table exponent.  Arbitrary changes of a maximizing
posterior only change the contents of this table.  No posterior
derivative or reciprocal probability denominator occurs.
\end{proof}

\begin{proof}[Proof of \Cref{prop:candidate-upgrade}]
Candidate existence means existence of a parent, node, chart and test;
the auxiliary construction proving it may use much sparser
configurations.  Only the residual is used to certify the candidate's
mass floor.

Use the preparations of
\Cref{prop:tangent-palette,a03:finite-state-conditioning}, retaining
\(1-o(1)\) of reference mass, including the pure caps and joint
end-history bounds with \(\pi_A\).  Bounded varying exponents use
neighboring widths in the increasing grids, at a \(K\) loss.
Aggregate with the fixed parent weights \(p_A\) and the unnormalized
submeasures of \(\lambda_A^0\).  Subsequent deletions change neither
reference law, so substantial mixture coverage is substantial
coverage on the reference path.

\paragraph{Coverage by one fixed candidate scenario.}
First fix a gain \(c>0\), a candidate type, and a bounded range
of path nodes in tangent units.  Use a fixed power slack
\(N^\sigma\) in the tests and mass threshold, where \(\sigma>0\)
will satisfy the bounds below.  In each chart \(Q\), choose
tests successively that carry the threshold mass on lines not
yet assigned.

Assign to each chosen test all still available \(Q\)-trajectories
satisfying its line inequalities.  For time candidates the value
and derivative upper bounds hold throughout the block, and the
derivative lower bound holds somewhere in it.  Only occurrences
where that lower bound holds count toward the mass threshold.
Eligibility may be measured under the unnormalized parent prior
restricted to prepared path incidences.

Disjointness makes this priority assignment stop after at most
\(N^{dc+\sigma}\) choices per \(Q\) for time candidates, or
\(N^\sigma\) for strips, with a \(K\) loss if path normalization
is used.  Run it separately for each node and scenario.
The retained geometric line predicates record the chosen tests;
they need not encode the analytical mass calculation that certified
their eligibility.


Choose a countable menu of scenarios.  Each fixes a type, a bounded
tangent-depth range, a positive gain \(c\), and a sufficiently small
power slack \(\sigma\); include arbitrarily fine choices of gains and
slacks.  The smallness conditions below are uniform on compact ranges
of positive gains.  If a residual candidate has gain \(c_n\to c_*>0\),
choose a menu gain sufficiently close to \(c_*\) and a slack \(\sigma\)
for which the difference of gains, multiplied by every fixed exponent
in the tests and mass floor, is less than \(\sigma/4\).  The candidate's
subpower losses and the depth-grid rounding use the remaining slack
for all sufficiently late stages.  In particular the change in the
factor \(N^{-dc}\) is included in this choice; it is not ignored.
Every candidate with positive limiting gain is thus eligible for
some fixed scenario eventually.

For each available scenario and node, make the finite greedy assignment
maximal.  The selections for different scenarios are made separately
on the same prepared path.  Suppose that the union of every fixed
finite menu has reference coverage tending to zero.  Choose
\(J_n\to\infty\) sufficiently slowly that the union of the first
\(J_n\) scenarios still has coverage tending to zero.  Remove
\emph{all} path incidences lying on their assigned trajectory pieces
throughout the respective blocks.  The remaining path has substantial
mass.  The candidate hypothesis gives, on a further subsequence,
a candidate of positive limiting gain in a bounded depth range.
The preceding menu choice supplies a fixed compatible scenario,
eventually among the first \(J_n\).  Its witnessing trajectories are
still free for that scenario in the relevant chart, since every
incidence on every earlier assigned piece was removed.  This contradicts
maximality.

Hence a fixed finite menu, of size \(J\), covers a fixed reference
fraction \(\eta_0>0\) on a subsequence.  Pigeonholing in that finite
menu fixes one scenario on a further subsequence.  Choosing one
available node per sample, among at most \(G_n\), leaves coverage
\begin{equation}
\beta_n\ge\frac{\eta_0}{J G_n}.
 \label{a03:covered-fraction}
\end{equation}
The node can vary with the sample; its depth remains in the fixed
scenario's bounded range.  By \Cref{a03:mass-diagonal},
\(\beta_n\log N_0\to\infty\).  Moreover \(\beta_n^{-1}=N^{o(1)}\)
by the diagonal preparation.  These are distinct facts: the latter
permits the analytical comparison, while the former supplies the
actual output mass.

The following comparisons lose only negligible fractions of this
coverage or bounded factors.  In the time case retain the derivative
lower bound \(N^{-2\sigma}\) at the covered incidences.  The derivative
on each assigned line is constant or affine, with maximum absolute
value at least \(N^{-\sigma}\), so the excluded times have relative
length \(O(N^{-\sigma})\).  Disjointness of assignments and
\Cref{a03:prepared-parent-law} bound their total mass by a fixed
negative power, hence by \(o(\beta_n)\).

\paragraph{Conversion of strip coverage.} Take a true label deeper
within the same parent, at a fixed target depth large enough that
its velocity-strip thickness is \(\ll gN^{-c}\); here \(\ell>0\).
Let \(A_{\rm deep}\) be its row, and put
\(D=R-A_{\rm deep}=(D_0,D_1)\).  On a counted incidence,
\[
 |D_0+D_1v|\le CgN^{-c+\sigma}.
\]
If \(\|D\|>gN^{-c/2}\), boundedness \(|v|\le K\) implies
either that no such incidence exists or that
\(|D_1|\ge gN^{-c/2}/K\), after enlarging \(K\).
The possible velocities then lie in an interval of length
\begin{equation}
 \frac{2CgN^{-c+\sigma}}{|D_1|}
       \le KN^{-c/2+\sigma}\le N^{-c/3}.
 \label{a03:strip-velocity-saving}
\end{equation}
For each deep history \(E\), its ancestor \(Q_s\) has at most
\(N^\sigma\) assigned rows.  The joint numerator bound
\Cref{a03:good-joint-numerator} and the angular ceiling therefore
charge all these misaligned rows by at most
\[
 K N^{\sigma-Sc/3}\lambda_{A,\mathrm{pre}}(E).
\]
Use fine velocity bins and bounded enlargements to cover the
intervals.  Summing over histories, labels, and parents with
weights \(p_A\) gives total mass at most
\(KN^{\sigma-Sc/3}\).  For example, choose
\(\sigma<\min(c/12,Sc/12)\).  This is a fixed-power saving and
hence is \(o(\beta_n)\), since \(\beta_n=N^{-o(1)}\).
Thus all but a negligible fraction of counted coverage satisfies
\begin{equation}
 \|R-A_{\rm deep}\|\le gN^{-c/2}.
 \label{a03:strip-row-alignment}
\end{equation}

The unbinned rows stored in the priority assignment also have
coefficients bounded by a fixed power of \(N\).  A row with much
larger norm and any eligible lines would confine bounded velocities
to an interval of arbitrarily small fixed-power length.  The same
joint bound, summed over the preceding histories in \(Q\), charges
that interval by its palette cost times \(Kq_s\nu_Q\).
Here the preceding history mass is at most \(Kq_s\nu_Q\) by
\Cref{a03:prepared-parent-law}.  For a sufficiently large fixed
power this contradicts the greedy eligibility threshold.

Set \(\tau=(J_s/q_s^2)N^{-c/4}\), using the actual relative
\(J_s,q_s\) at \(Q\), and bin the two coefficients of \(R\)
at width \(\tau\), with values \(R_0^b,R_1^b\).
On aligned events the deeper path label gives an actual list
for this bin at original \(p_L\), since
\(gN^{-c/2}\ll\tau\).  The row \(R\) itself then has coefficients
bounded by \(K_{\rm path}^{O(1)}\), using the actual deep-row
bounds from \(Z\)-projection fullness.  Omit predictors violating
these bounds from the output.  On assigned lines,
\begin{equation}
|Y'-R_0^b-R_1^b Z'|\le K_{\rm path}^{O(1)}\tau
\end{equation}
by the actual bound on \(v=Z'\) and the positive power buffer.
Also bin \(Y-R_1^bZ\) at the block midpoint at width \(q_s\tau\).
Given the observed positions and velocity-row bin, this position
bin has actually subpower ambiguity.  The output label consists of the old label \(Q\), its binned row, and
this position bin.  Retain each trajectory's earlier priority assignment and merge
assignments with the same three entries.  This assignment is fixed throughout the
block.  The row and position bins, rather than the predictor's identity,
are the data to be recovered from the original observation.

Keep the old \(Z\)-center.  For \(Y-R_1^bZ\), use the selected
partition value moving with derivative \(R_0^b\), and take spatial
columns \(q_s(1,R_1^b)\) and \(q_s(0,\tau)\) in parent
\((Z,Y)\) coordinates.  This geometry has actual path-slack bounds
throughout the block.  Its determinant in parent coordinates is
\(q_s^2\tau=J_sN^{-c/4}\), so composition through the parent
multiplies the old cumulative determinant at \(Q\) by \(N^{-c/4}\).
Keep the true test, depth, and time.  The resulting atlas has the
forbidden narrowness improvement and the required aggregate list
success on the original path; its tempered implementation is given
below.

\paragraph{Comparison of the time candidates.}
We compare candidate graphs compatible with the same original
observation.  They agree at its base point to the candidate fitting
accuracy.  We will prove that, except
on a negligible fraction of compatible pairs, they are also close on
a whole short comparison box.  Their finite jets can then replace the
candidate indices as listable chart labels.

Work in each assigned \(Q_s\), using its actual normalized base coordinates \(u\) (unit-block time and spatial coordinates via its moving affine center and frame); \(w\) remains in parent units. Write \(a_Q\) for the actual node thickness in these parent units, \(a_Q/a\sim_{\log,N}1\). Any subpower factors in the following preparatory comparisons can be absorbed using \(N^\sigma\), for fixed \(\sigma\) and sufficiently late stages. The constants in exponents \(O(\sigma)\) below can be fixed independently of \(\sigma\) small. Write \(m_Q=p_A q_s\nu_Q\), so \(\sum_Q m_Q\le1\) summing over parents and blocks at the chosen node.

Each index (assigned test)'s base marginal on the counted occurrences costs at most
\begin{equation}
m_Q N^{-dc+O(\sigma)}\,du.                                                     \label{a03:candidate-marginal}
\end{equation}
Indeed the candidate value and lower derivative tests require only \(O(N^{O(\sigma)})\) hidden bins of width \(a'\) at each exact base (use monotone intervals of the quadratic); use the \(K(a')^d\) cap and Jacobian \(q_s J_s\), with \(\nu_Q\sim_{\log,N}a^d J_s\). This bound holds also on each test's original matched incidence set witnessing the greedy mass threshold. That set's projection in \(u\) has volume \(\ge N^{-O(\sigma)}\), by the threshold and \Cref{a03:candidate-marginal} (likewise if using the unnormalized parent prior for eligibility).

Convert each assigned test and the true \(Q_s\) test into local
graphs in \(w\).  Write the test as \(P\), with
\(a_{\rm fit}=a'\) for a candidate and \(a_{\rm fit}=a\) for
the true test.  At late stages its whole-block value bound is
\(a_{\rm fit}N^\sigma\), its whole-block derivative upper bound
is \(N^\sigma\), and its derivative lower bound at counted points
is \(N^{-2\sigma}\), up to harmless fixed constants.
If \(|P_{ww}|\ge N^{-C\sigma}/a_{\rm fit}\), for a sufficiently
large fixed \(C\), say \(20\), use the affine quadratic center.
It stays within \(a_{\rm fit}N^{O(\sigma)}\) of the entire
assigned line piece; if an equation is needed, use \(w\) minus
this center.

For the remaining tests the discriminant
\begin{equation}
\mathcal D=P_w^2-2P_{ww}P
\end{equation}
is independent of \(w\) and at counted points positive and comparable to \(P_w^2\). On complex enlargements of the entire normalized chart box (\(K_{\rm path}\)-sized ranges in \(u\) allowed) this polynomial and its first derivatives are bounded by \(N^{O(\sigma)}\). Indeed its upper bound on the matched projection just discussed follows immediately from the tests there, for both the assigned candidate and true test; extend by polynomial Remez. This uses the initial matching before greedy enlargement for each assigned candidate, not a new floor on the path in each subbox.

Subdivide the time of \(Q_s\) dyadically at relative size \(\rho\)
comparable to \(N^{-C'\sigma}\), with \(C'\) sufficiently large and
fixed.  In each short block, bin every line's normalized spatial
position at the midpoint into \(\rho\)-boxes.  There are at most
\(N^{O(\sigma)}\) resulting base subboxes, including time, per \(Q\).
Given \(Q\) and the exact base position, these subboxes have lists
at the original \(p_L\) of size bounded by actual path-slack powers:
motion in normalized coordinates over a short block is at most
\(K_{\rm path}^{O(1)}\rho\).

Around each subbox center, take comparison balls or polydiscs with
radius a fixed multiple of \(K_{\rm path}^{O(1)}\rho\), large enough
to contain the corresponding real line pieces with a fixed extra
holomorphic margin.  Keep only test--subbox choices certified by a
counted occurrence, which supplies the derivative lower bound at one
point.  The chart derivative bound for \(\mathcal D\), together with
the choice of \(C'\), makes its variation throughout this comparison
domain much smaller than \(N^{-4\sigma}\).  For each nonreplaced
quadratic test, the discriminant therefore stays nonzero and gives
simple holomorphic branches throughout the domain, real on the real
box.  For a linear equation in \(w\), the same check keeps its
coefficient nonzero.

For such nonreplaced tests, on any assigned line for these choices the low-curvature cutoff guarantees \(2|P_{ww}P|\ll\mathcal D\) over the whole real short piece, so the sign of \(P_w\) is constant there, picking one branch \(f\). At every time \(w\) has distance at most \(O(a_{\rm fit}N^{O(\sigma)})\) from that branch: \(P_w\) has the same sign there and at \(w\), comparable in absolute value to \(\sqrt{\mathcal D}\), so the derivative lower bound works on the intervening interval. Replaced tests simply use their affine centers. Thus branch signs for true and candidate tests can be assigned invariantly over the short block. Denote the true comparison graph in one such choice by \(f_0\).

Let \(\mu_{\mathrm{cov}}\) be the unnormalized law of counted covered
occurrences, including the parent weights.  Set
\(O=(p_L,Q_s,\text{subbox},\text{true branch})\), with history
included in \(Q_s\), and let \(\kappa_O\) be the conditional
probability of an occurrence under this law.  Sample \(O\) with
its covered-law pushforward, then sample two independent posteriors:
\begin{equation}
 \mu_{\mathrm{pair}}
   =\int \kappa_O\otimes\kappa_O\,d(O_*\mu_{\mathrm{cov}})(O).
 \label{a03:covered-pair-law}
\end{equation}
Each marginal is \(\mu_{\mathrm{cov}}\), and the pair law retains
its unnormalized total mass, at least \(N^{-o(1)}\).
The candidate graphs agree at the sampled common base point to
\(a'N^{O(\sigma)}\), because original \(p_L\), given the parent,
determines \(w\) much more finely.  In version 1 the subtracted
parent graph is fixed by the chart and exact base, so the same
conclusion holds.

Fix \(d_0=(1+d)/2<1\).  Let \(\eps_*\) be the tolerance in
\Cref{a02:local-graph-matching} for these fixed comparison domains
and degrees.  Decrease it to at most one, and set
\begin{equation}
 \eps=\eps_*/2,\qquad \theta=\eps_*/8.
 \label{a03:matching-tolerances}
\end{equation}
We will choose \(\sigma\) after these constants.  We claim that,
outside an \(o(1)\) relative fraction of the pair law, the candidate
graphs agree in norm on the comparison real box, with its fixed
margin around the used pieces, to \(aN^{-2\theta c}\).

First discard entries (index, subbox, candidate and true branch
signs) with marginal mass below \(m_QN^{-dc-C''\sigma}\).
The index and subbox counts make their total mass negligible on
either marginal when \(C''\) is sufficiently large.
For each remaining entry, the candidate graph differs from
\(f_0\) by at most \(aN^{O(\sigma)}\) on its incidence base
projection.  By \Cref{a03:candidate-marginal}, this projection
has relative measure at least \(N^{-O(\sigma)}\) in the comparison
box.  Apply \Cref{lem:algebraic-norm} on the interiors with fixed
margins to obtain the same norm bound, with an adjusted
\(O(\sigma)\) exponent.

Suppose a nonvanishing relative fraction of the pair law were bad.
The light-entry omissions leave a group
\((Q,\text{subbox},\text{true branch})\) of bad pairs of mass
at least \(m_QN^{-O(\sigma)}\); the weighted count justifying
this choice is given below.  Normalize that group's bad pairs
to probability.  We compare its two candidate lists by
\Cref{a02:local-graph-matching}, relative to the common true graph
\(f_0\), with \(M=a\) and \(\Delta=a'\).
After rescaling the base, \Cref{a03:candidate-marginal} gives
index-base ceilings with weights \(N^{-dc}\) and losses
\(N^{O(\sigma)}\).  To check that they remain valid for probability
weights in each color, we also account for the bad-pair normalization.
Indeed, if \(W_c\) is the sum of \(N^{-dc}\) over that color's retained
indices and \(B\) is the unnormalized bad-pair mass in the selected group,
the index density bound becomes
\begin{equation}
\frac{f_i}{B}
 \le N^{O(\sigma)}
       \frac{m_Q W_c}{B}\,
       \frac{N^{-dc}}{W_c}.
\end{equation}
The list count gives \(W_c\le N^{O(\sigma)}\), and the group choice gives
\(B\ge m_QN^{-O(\sigma)}\).  Thus the multiplier of the probability
weight \(N^{-dc}/W_c\) is \(N^{O(\sigma)}\), as required.
The existence of such a group uses the sum
\(\sum_Qm_Q\le1\) and the polynomial-in-\(N^\sigma\) subbox count.
If a nonvanishing \emph{relative} fraction of the covered pair law were
bad, its absolute mass would be at least a fixed multiple of
\(\beta_n=N^{-o(1)}\).  Light-entry errors are power-small and hence
\(o(\beta_n)\); the factor \(\beta_n^{-1}\) is absorbed into the
fixed-power comparison allowance.

Here is the tolerance check.  Put \(R=M/\Delta=N^c\), with
\(M=a\) and \(\Delta=a'\).  All functions are holomorphic sheets
with bounded-degree relations.  A bad pair has norm gap at least
\[
 aN^{-2\theta c}=MR^{-2\theta}>MR^{-\eps}.
\]
The list counts are at most \(R^{d+C\sigma/c+o(1)}\).
The norm upper bounds relative to the true graph \(f_0\), the
pointwise matching precision, and the normalized marginal multipliers
just computed have losses at most \(R^{C\sigma/c+o(1)}\), for
finitely many fixed constants \(C\).  Choose \(\sigma\) so that
every \(C\sigma/c<\eps_*/8\).  At late stages the remaining
\(o(1)\) losses, including \(\beta_n^{-1}\), fit within another
\(\eps_*/8\).  The normalized bad-pair law has mass one, so it
also satisfies the required lower mass bound.  Since \(d<d_0\),
all hypotheses of \Cref{a02:local-graph-matching} hold with
\(\eps=\eps_*/2\), a contradiction.  This proves the claim.

\paragraph{The improved time charts.}
Bin sufficiently many candidate branch jets at the subbox center (in comparison-ball scaled coordinates) to steps \(aN^{-\theta c}\). Pairs with the proved closeness give bounded neighbor bins by the norm rule. For clarity, if \(\pi_O\) is the posterior law of the actual jet bin
at a conditioning observation \(O\), and \(\mathcal N(b)\) is the
bounded neighbor list of a bin \(b\), then
\begin{equation}
\sum_b\pi_O(b)\pi_O(\mathcal N(b))
 =
 \Pp\{B_2\in\mathcal N(B_1)\mid O\}.
 \label{a03:posterior-neighbors}
\end{equation}
Choose a maximizing \(b\) separately for each \(O\).  Its neighbor list
has success at least the displayed average.  Integrating preserves
\(1-o(1)\) of the covered mass once the bad-pair fraction is \(o(1)\);
in particular it preserves a fixed positive fraction.  This argument retains aggregate covered mass; it does not select an
individual candidate that may carry only inverse-subpower mass.

Define the preliminary geometric label
\[
 c_{\mathrm{geom}}=(Q_s,\text{short block},\text{spatial subbox},
             \text{true branch sign},\text{candidate jet bin}).
\]
Keep the earlier priority assignment and its invariant branch signs,
and merge assignments with the same \(c_{\mathrm{geom}}\).  This gives a trajectory
assignment fixed throughout the short block.  The original candidate
index is not part of \(c_{\mathrm{geom}}\).  The conditioning entries preceding the
jet bin have only actual path-slack ambiguity at the original
\(p_L\), so the posterior-neighbor lists give lists for this full
preliminary label with the same type of bound.

For each nonempty jet bin in its group, choose one representative
equation \(\widehat P\) and its assigned branch \(\hat f\).
The representative is a fixed function of \(c_{\mathrm{geom}}\).  Every member
branch in this bin differs from \(\hat f\) by
\(O(aN^{-\theta c})\) on the real box: sufficiently many center
jets control the norm of their holomorphic algebraic difference.
Thus
\[
 |w-\hat f|\lesssim aN^{-\theta c}
\]
on each member's entire short line piece, after choosing \(\sigma\)
small enough in the graph approximations.  We now turn the
representative into an admissible chart test.

\begin{itemize}
\item If \(|\widehat P_{ww}|\le (aN^{-\theta c/2})^{-1}\), normalize by the absolute \(w\)-derivative on the branch at box center (this equals 1 for explicit tests, notably in version 1). In the simple-root case this factor is \(\ge N^{-O(\sigma)}\) and comparable uniformly, to constants, to the absolute branch derivative throughout by discriminant stability. Passing from branch to member \(w\) changes the derivative by \(O(N^{-\theta c/2})\). Thus the normalized equation has actual order-one derivative bounds, value \(\lesssim aN^{-\theta c}\) and curvature \(\le a^{-1}N^{\theta c/2+O(\sigma)}\) on the required data. These fit at thickness \(a_Q N^{-2\theta c/3}\).
\item For larger curvature, the representative branch is within \(aN^{-\theta c/2+O(\sigma)}\) of its affine center by the branch derivative upper bound. Use \(w\) minus that affine center as test with thickness \(a_Q N^{-\theta c/3}\).

\end{itemize}
Choose \(\sigma\) after \(\theta\) to meet the preceding matching
tolerances and graph-approximation bounds.  Also impose
\(C'\sigma<k\theta c/12\).  The strict buffers absorb the
subpower discrepancy \(a_Q/a\).  Select a representative type
carrying a fixed fraction of the list success, so the output
has a common depth.

Use the true spatial frame scaled by \(\rho\), with center
adjusted to the subbox's midpoint spatial bin in moving chart
coordinates.  Its bounds hold throughout with actual path slacks.
Restore original coordinates by \Cref{lem:normalization}:
thickness and tests acquire the parent-depth factor, with derivative
units restored.  In version 2, \(w\) uses the resulting parent
\(B\)-scale; in version 1 the parent graph subtraction preserves
the special test form.  The path labels, subbox, true branch sign,
and deterministic bin lists together cost only actual path-slack
and logarithmic powers.

Keep the preliminary label \(c_{\mathrm{geom}}\); its chosen representative already
determines the curvature type and the final test.  Retain only assigned
trajectories satisfying the resulting whole-block geometric, value,
and derivative upper bounds, and impose the derivative lower bound on
selected incidences.  The comparisons above show that the retained
list successes satisfy these tests.  We have therefore constructed
geometric charts and full-label lists on the covered analytical law.

The depth gain is at least \((\theta c/3)\epsilon_n\), whereas
the additional time exponent is
\(C'\sigma\epsilon_n+o(\epsilon_n)\).
Our choice of \(\sigma\) makes this less than half \(k\) times
the depth gain.  Either representative type therefore gives the
forbidden improved-depth atlas.

\paragraph{Finite encoding and return to the original law.}
Both constructors have now specified their chart geometry, tests, full
labels, and block-invariant trajectory assignments.  We verify that these
objects have bounded original-coordinate complexity and recover their
lists on the original tempered path.

Store the ordered quadratic tests or row predictors at the single chosen
node.  The assignments use membership in the old chart's trajectory set,
whole-block line inequalities, and finite priority rules.  Subbox
certification flags, branch signs, merger tables, and representative tests
are fixed data for each chart and test; a simple branch is selected by
its midpoint derivative sign.  Nonempty jet bins may be reindexed by
discrete identifiers, so numerical jets need not enter a predicate.
The analytical masks and the posterior search used to choose the lists
are not part of these geometric assignments.

\paragraph{Coefficient bounds.} Here is an explicit coefficient bound for every time candidate.
In normalized chart coordinates write
\(P(u,w)=A w^2+B(u)w+C(u)\), with \(B\) affine and \(C\) quadratic.
Curvature gives \(|A|\le N^{s+c+O(\sigma)}\).  The original threshold
witness, before greedy enlargement, has base projection of volume
\(\delta\ge N^{-C_{\mathcal S}\sigma}\) by
\Cref{a03:candidate-marginal}; its chart weight cancels when the floor is
divided by the ceiling.  On that witness
\(|w|\le W\le N_0^{P_0(M_{\mathrm{path}})}\) by original problem and
path bounds.  The derivative and value estimates yield
\begin{equation}
|B(u)|\le N^\sigma+2|A|W,\qquad
 |C(u)|\le a'N^\sigma+|A|W^2+|B(u)|W .
 \label{a03:coefficient-recovery}
\end{equation}
Fixed-degree polynomial Remez from this projection costs at most
\(\delta^{-C}\); it bounds all coefficients of \(B,C\) by
\(N_0^{P_{\mathcal S}(M_{\mathrm{path}})}\).
This uses a bound for \(w\) in original exponent units and does not
assume that its affine part is bounded in tangent units. Inverse curvatures for affine centers are needed only above cutoffs; rescaling by a branch derivative uses the stability lower bound. Raw strip coefficients were bounded above. Coordinate compositions back to the original world are all polynomially controlled. Counts of new choices per path chart cost only fixed tangent powers. The tests for an entire line/block here have fixed-degree semialgebraic complexity (if quantifying over block time do that only for each individual test separately).
\paragraph{Actual lists.} On the masked successes the full output labels have lists at the
original \(p_L\) of size at most
\((K_{\mathrm{path}}\log N_0)^{r_{\mathcal S}}\), with
\(r_{\mathcal S}\) fixed by the scenario.  Apply
\Cref{a03:actual-lists} to the original unnormalized reference path
and its masked covered submeasure, with the geometric assignment just
defined.  This produces a tempered cell-list selection with the same
list size and only \(O(N_0^{-1})\) original mass loss.  In particular,
one does not flatten the analytical masks or encode their posterior
rules.  The final selector observes only the original exact base and
original discrete observation entries.
\paragraph{Complexity and mass.}
The coefficient and predicate bounds above, followed by
\Cref{a03:actual-lists}, bound every original-coordinate complexity
exponent by a polynomial in \(M_{\mathrm{path}}\) depending only on
the fixed scenario.  This includes the flattening depth and list-table
count.  Only the chosen scenario is encoded.  Its polynomial bound
fits \(M^\#\) at all sufficiently large stages by
\Cref{a03:complexity-guard}; its fixed power of actual path factors
and logarithms likewise fits the allowance \(n\).

The comparison and representative choice retain some fixed fraction
\(c_0>0\) of \(\beta_n\).  After original-law list recovery, the
original final success is therefore at least
\begin{equation}
\frac{c_0\eta_0}{J G_n}\,\mathfrak p_{\mathrm{path}}
       -O(N_0^{-1})
 >
 \frac{\mathfrak p_{\mathrm{path}}}{\log N_0}
 \label{a03:final-coverage}
\end{equation}
for large stages, by \Cref{a03:mass-diagonal}.  The original mass loss
is negligible since
\(N_0^{-1}\log N_0/\mathfrak p_{\mathrm{path}}
 \le N_0^{-1}K_{\mathrm{path}}\log N_0\to0\).
The fixed positive tangent gain exceeds \(\gamma_n\), and its time
cost is less than half \(k\) times its gain.  All requirements of
\Cref{a03:finite-guard} are now satisfied, giving the contradiction.

\end{proof}



\subsection{Velocity breadth in isotropic parents}

The isotropic argument will use finite point states to control gradients
of packet graphs.  A large gradient would predict the velocity in a
power-thin affine strip.  We now rule out such concentration on a
substantial family by constructing a narrower parent atlas.  The
original-law list recovery and finite encoding just established realize
this width improvement while keeping the parent depth and horizon fixed.

\begin{corollary}[Breadth before sparse selections]\label{a03:breadth}
In version 2 with \(\ell=0\), let a substantial path residual carry a
measurable discrete state, including its tangent parent.  Suppose the
state is known on these events at the original \(p_L\) with list size
bounded by a fixed power of actual path factors and logarithms, and has
at most \(K\) possible values given the used deep end histories.
For every fixed \(\kappa,u>0\), the total residual probability of states
whose conditional velocity law puts mass at least \(\kappa\) within
distance \(N^{-u}\) of an affine line tends to zero.

In version 1 the corresponding assertion for intervals follows from
the angular cap.  These statements can be prepared at reciprocal-subpower
separation thresholds before arbitrary sparse selections.  They persist
under selections of whole conditioning states, or controlled
fixed-fraction selections within them; no version-2 breadth assertion
is made for every velocity-dependent inverse-subpower subfamily.
\end{corollary}
\begin{proof}
If the version-2 assertion fails, on a subsequence there are
state-predicted strips of the indicated width carrying a fixed positive
amount of residual mass, still on substantial events of the prepared
good law.
\paragraph{A finite family of heavy strips.}
Use unit normal and offset pairs \((n_v,b_v')\) for those predicted lines, with Euclidean parameter distance modulo simultaneous sign. Parent velocities are bounded by \(T\ge1\) of order actual path-slack powers; offsets needed then cost \(O(T)\). On almost all the narrow events the predicted strip, say enlarged to \(2N^{-u}\) using finer velocity bins, has parent palette \(\pi\)-weight at least \(\kappa_N=1/K\) for sufficiently small reciprocal-subpower threshold. Indeed those of smaller palette mass cost negligibly by summing the unnormalized joint bound of velocity with the end history using the subpower list there (choose \(\kappa_N\) small relative to all those losses). State predictions need only be specified on occurring states; their choices are finite here.

Per parent take maximal parameter-separated representatives at spacing \(r_1=N^{-u/2}\) for such heavy strips. Their number \(M_1\) is \(\le K\). For two separated ones with common narrow-band points in \(|v|\le T\), the angle sine is at least a constant times \(r_1/(1+T)\): take the sign with closer normals; offset difference there is at most \(4N^{-u}+T|n_{v,1}-n_{v,2}|\). Thus intersection costs a ball of diameter \(O((1+T)N^{-u}/r_1)\le N^{-u/3}\), hence a fixed angular power in \(\pi\). By Cauchy–Schwarz on the sum of strip indicators,
\begin{equation}
(M_1\kappa_N)^2\le M_1+K N^{-Su/3} M_1^2,
\end{equation}
giving the claim. Moreover for representatives separated by \(>r_2=N^{-u/4}\), their \textbf{enlarged} strips of widths \(10(1+T)r_1\) still have intersections of negligible palette mass even summed over all pairs: the same reasoning gives diameter \(\lesssim (1+T)^2r_1/r_2\) or no common point in the used range.

\paragraph{Trajectory assignments and the narrower atlas.}
On narrow events the velocity therefore lies in an enlarged representative strip, and any choice of such containing strip agrees to \(O(r_2)\) in parameter (modulo sign) with the predicted pair outside negligible mass. Indeed there is one within \(r_1\) of the prediction by maximality, and the excluded overlaps depend just on parent and velocity and cost little by the palette bound. Assign the first containing representative by \textbf{velocity alone} on parent trajectories, invariantly over the block. Bin its parameters (with one orientation each) to rectangular \(r_2\)-bins; this costs only boundedly many options given the state prediction on successful events. Fix one unit normal and offset pair \((\hat n,\hat b)\) per occurring bin consistent with it. Then
\begin{equation}
|\hat n\cdot y'-\hat b|\lesssim (1+T)r_2
\end{equation}
on assigned pieces. Bin also the corresponding \(\hat n\)-position at parent-block midpoint to width \(r_2\) in parent units, known with actual slack up to lists via observed base and this velocity prediction. Merge assignments using those two bin labels. Use a frame with unit tangential column and width \(r_2\) along \(\hat n\) in parent spatial coordinates (move the transverse bin with derivative \(\hat b\)); compose with the parent frame, keeping parent test, depth and horizon. All slacks thus cost only actual path bounds to fixed powers, and the
determinant has gained the new narrowness factor.

On the successful incidences, the assumed state knowledge supplies
actual lists for the full output labels at the original \(p_L\).
Store the finite representative and bin assignments together with the
original path predicates.  Apply \Cref{a03:actual-lists} to this
geometric label map on the original path and its successful submeasure.
The analytical list success is then realized by tempered tables with
the same list size and only \(O(N_0^{-1})\) original mass loss; the
state predictions themselves need not be encoded.  As in the preceding construction, the original mass and
fixed-complexity bounds fit \Cref{a03:finite-guard}.  This contradicts
the width improvement excluded at the parent.

In version 1 the corresponding avoidance of narrow velocity intervals on substantial families already follows from the angular cap and the deep-history joint bound (summing the subpower state ambiguities there). These almost-all breadth/separation conclusions can use reciprocal-subpower accuracies by taking fixed-power thresholds to zero sufficiently slowly. In applications requiring general position, one can prepare them before sparse selections and keep them through selections of whole conditioning states (or with small fixed-fraction losses controlled there). No breadth in version 2 on every \(K^{-1}\) velocity-dependent selection is claimed.

\end{proof}
